\documentclass[12pt]{article}


\usepackage[T2A]{fontenc}
\usepackage[utf8]{inputenc}
\usepackage[serbianc]{babel} 

\usepackage{graphicx} 
\usepackage{hyperref}
\usepackage{xcolor} 
\usepackage{amsmath, amsthm} 
\usepackage{wrapfig} 


\hypersetup{
    colorlinks=true,   
    linkcolor=blue,    
    urlcolor=cyan,      
}


\newtheorem{definicija}{Дефиниција}

\begin{document}



\begin{titlepage}
    \center 
    \textsc{\LARGE Математички факултет}\\[2cm] 
    \textsc{\Large Семинарски рад}\\[0.5cm]
   
    \vspace{1cm}
    {\LARGE\bfseries Дигитални близанци (Digital Twins) \\[0.4cm] \large Виртуелне копије физичких система} \\[1cm]
    
    \vspace{4cm}
    \begin{minipage}{0.4\textwidth}
        \begin{flushleft} \large
            \textit{Студент:}\\
            Ђорђе Шљивић 146/2025
        \end{flushleft}
    \end{minipage}
    \begin{minipage}{0.4\textwidth}
        \begin{flushright} \large
            \textit{Професор:}\\
            Данијела Симић
 \begin{figure}
            \centering
          
            
            \label{fig:placeholder}
        \end{figure}
                \end{flushright}
    \end{minipage}
    
    \vfill
    {\large Београд, \today} 
\end{titlepage}

\tableofcontents
\newpage


\section{Увод}
Дигитални близанци (\textit{Digital Twins}) представљају виртуелне репрезентације реалних објеката. Овај концепт омогућава симулацију понашања система у виртуелном окружењу.

\begin{definicija}
Дигитални близанац је интегрисана виртуелна репрезентација стварног ентитета која је синхронизована са физичким системом у реалном времену.
\end{definicija}


\section{Карактеристике и модели}

\subsection{Математички модел и верност}
Ефикасност близанца често зависи од функције верности $f(x)$, где се грешка симулације може представити као:
$$ \epsilon = \sum_{i=1}^{n} (y_i - \hat{y}_i)^2 $$
где је $y_i$ реални податак, а $\hat{y}_i$ вредност из виртуелног модела. Такође, користимо познате релације попут енергетске ефикасности $E = mc^2$ у термодинамичким симулацијама.

\subsection{Визуелизација система}

\begin{wrapfigure}{l}{0.4\textwidth}
    \centering
    \includegraphics[width=0.35\textwidth]{digital-twin-component-view(1).png} 
\begin{figure}
        \centering
        \includegraphics[width=0.5\linewidth]{digital-twin-component-view(1).png}
        \caption{Enter Caption}
        \label{fig:placeholder}
    \end{figure}
        \caption{Шема симулације}
\end{wrapfigure}
Овде видимо како се подаци прикупљени преко сензора шаљу у базу. \textbf{\textcolor{blue}{Плавом бојом}} су означени сигурни преноси, док је \textbf{\textcolor{red}{црвеном}} означен ризик од кашњења. \textit{Кључно је постићи ниску латенцију} како би модел био ажуран.


\section{Области примене}
Дигитални близанци се користе у следећим секторима:
\begin{itemize}
    \item \textbf{Индустрија:} Оптимизација линија.
    \item \textbf{Енергетика:} Контрола мреже.
    \item \textbf{Медицина:} Моделовање органа.
\end{itemize}

\begin{table}[h]
\centering
\caption{Поређење традиционалне и дигиталне симулације}
\begin{tabular}{|l|c|c|}
\hline
\textbf{Параметар} & \textbf{Традиционално} & \textbf{Дигитални близанац} \\ \hline
Ажурирање & Статично & Реално време \\ \hline
Повратна веза & Недостаје & Дво smerна \\ \hline
Прецизност & Средња & Висока \\ \hline
\end{tabular}
\end{table}

\section{Закључак}
Можемо закључити да су дигитални близанци кључна технологија будућности. Иако је имплементација сложена, користи у виду \textbf{смањења трошкова} и \textit{веће безбедности} су незаменљиве.

\end{document}